\documentclass[10pt,a4paper]{amsart}
\usepackage[margin=19mm]{geometry}
\usepackage{amsmath,amssymb,mathtools}
\usepackage[scr=rsfs,cal=euler]{mathalfa}
\usepackage{xfrac}
\usepackage[unicode,hidelinks,bookmarks=false]{hyperref}
\newcommand{\ie}{i.e.\ }
\newcommand{\cf}{cf.\ }
\newcommand{\resp}{respectively\ }
\newcommand{\Subsection}{\subsection}
\providecommand{\nobreakdash}{\nobreak\hbox{-}}
\setlength{\emergencystretch}{2em}
\multlinegap0pt
\usepackage{amssymb}
\usepackage{float}
\usepackage{tikz}
\usepackage{enumerate}
\usepackage[shortlabels]{enumitem}
\usepackage{algorithm}\usepackage{algpseudocode}
\newtheorem{theorem}{Theorem}
\renewcommand{\c}[1]{\mathcal{#1}}
\newcommand{\m}[1]{\mathbf{#1}}
\DeclareMathOperator{\tr}{tr}
\title{{\large Spectral Shrinkage of Gaussian Entropic Optimal Transport}}
\author{Ho Yun}
\date{Source: arXiv:2512.19457v2 (CC BY 4.0)}
\begin{document}
\maketitle

\noindent\emph{Source and licence.} {\large Spectral Shrinkage of Gaussian Entropic Optimal Transport}. Version arXiv:2512.19457v2 (CC BY 4.0), distributed by arXiv under the Creative Commons Attribution 4.0 International licence, http://creativecommons.org/licenses/by/4.0/, as recorded in the arXiv licence metadata for this exact version and shown on the abstract page https://arxiv.org/abs/2512.19457v2. Full licence notice: \texttt{licenses/arxiv-2512.19457-CC-BY-4.0.txt}.\par\smallskip

\noindent\textit{Editorial scope.} Counted unit: the article's theorem thm:W2:exact (source0.tex lines 1050--1065). The article's complete original proof of it is delivered with it (source0.tex lines 1066--1118, one \texttt{proof} environment of 53 lines). No mathematics is written, repaired or re-derived: the statement and the proof are the article's own lines, in the article's own notation, and no retained line is substituted or re-flowed.  No further source-local context is needed: every cross-reference of every retained line resolves inside the two delivered spans.  Nothing was omitted from any retained span, so the package carries no omission record.  No retained line contains a citation, so the wrapper carries no bibliography and no bibliography entry.  No retained line contains a figure environment, an image-inclusion command or drawing code, so no image is delivered and the package declares no asset dependency.  Editorial formatting only, in addition: an AMS \texttt{amsart} preamble that restates the article's own declarations and macros used by the retained lines, namely the environments \texttt{theorem} on separate counters, together with the standard packages those lines need.  The retained environments print in this document as Theorem 1 in order of appearance, whereas the article prints the same items with the numbering of its own full document.  The complete native source of the article as distributed in the arXiv e-print for this exact version is bundled unchanged as \texttt{sources/191479/source0.tex}.
\begin{theorem}\label{thm:W2:exact}
Let $\varepsilon > 0$ and $\m{A}, \m{B} \in \c{B}_{1}^{+}(\c{H})$ with $\m{A} \leadsto \m{B}$. Then, the squared Wasserstein distance between the Gaussian EOT coupling $\pi_{\varepsilon} = \c{N}(\m{0}, \m{\Sigma}_{\varepsilon})$ and the canonical OT coupling $\pi_{0} = \c{N}(\m{0}, \m{\Sigma}_{0})$ is given by
\begin{align*}
    \c{W}_{2}^{2}(\pi_{\varepsilon}, \pi_{0}) = 2 \tr [\m{A} + \m{B}- \sqrt{\m{Q}_{\varepsilon}}],
\end{align*}
where
\begin{align*}
    \m{Q}_{\varepsilon} :&= (\m{G}_{0}^{*} \m{G}_{0})^{2} + (\m{M}_{0}^{*} \m{M}_{0})^{2} + (\m{G}_{0}^{*} \m{G}_{0}) \m{R}_{\varepsilon} (\m{M}_{0}^{*} \m{M}_{0}) + (\m{M}_{0}^{*} \m{M}_{0}) \m{R}_{\varepsilon} (\m{G}_{0}^{*} \m{G}_{0}). 
    % \\
    % &= \begin{pmatrix}
    % \m{G}_{0}^{*} \m{G}_{0} & \m{M}_{0}^{*} \m{M}_{0}
    % \end{pmatrix} \begin{pmatrix} \m{I} & \m{R}_{\varepsilon} \\ \m{R}_{\varepsilon} & \m{I} \end{pmatrix} \begin{pmatrix}
    % \m{G}_{0}^{*} \m{G}_{0} \\ \m{M}_{0}^{*} \m{M}_{0}
    % \end{pmatrix} \in \c{B}_{0}^{+} (\c{H}).
\end{align*}
\end{theorem}

\begin{proof}
Noting that
\begin{align*}
    \m{\Sigma}_{0} &= \left(
    \begin{array}{c|c}
    \m{A} & \m{G}_{0} \m{M}_{0}^{*} \\
    \hline
    \m{M}_{0} \m{G}_{0}^{*} & \m{B}
    \end{array}
    \right) = \left(
    \begin{array}{c}
    \m{G}_{0} \\
    \hline
    \m{M}_{0}
    \end{array}
    \right) \left(
    \begin{array}{c|c}
    \m{G}_{0}^{*} & \m{M}_{0}^{*}
    \end{array}
    \right),
\end{align*}
the spectrums of $\m{\Sigma}_{\varepsilon}^{1/2} \m{\Sigma}_{0} \m{\Sigma}_{\varepsilon}^{1/2}$ and
\begin{align*}
    \m{Q}_{\varepsilon} &= \left[\m{\Sigma}_{\varepsilon}^{1/2} \left(
    \begin{array}{c}
    \m{G}_{0} \\
    \hline
    \m{M}_{0}
    \end{array}
    \right) \right]^{*} \left[\m{\Sigma}_{\varepsilon}^{1/2} \left(
    \begin{array}{c}
    \m{G}_{0} \\
    \hline
    \m{M}_{0}
    \end{array}
    \right) \right] = \left(
    \begin{array}{c|c}
    \m{G}_{0}^{*} & \m{M}_{0}^{*}
    \end{array}
    \right) \m{\Sigma}_{\varepsilon} \left(
    \begin{array}{c}
    \m{G}_{0} \\
    \hline
    \m{M}_{0}
    \end{array}
    \right) \\
    &= (\m{G}_{0}^{*} \m{G}_{0})^{2} + (\m{M}_{0}^{*} \m{M}_{0})^{2} + (\m{G}_{0}^{*} \m{G}_{0}) \m{R}_{\varepsilon} (\m{M}_{0}^{*} \m{M}_{0}) + (\m{M}^{*} \m{M}) \m{R}_{\varepsilon} (\m{G}_{0}^{*} \m{G}_{0}),
\end{align*}
coincide. Therefore, we conclude that
\begin{align*}
    \c{W}_{2}^{2}(\pi_{\varepsilon}, \pi_{0}) = \tr[\m{\Sigma}_{\varepsilon}] + \tr[\m{\Sigma}_{0}] - 2 \tr [\m{\Sigma}_{\varepsilon}^{1/2} \m{\Sigma}_{0} \m{\Sigma}_{\varepsilon}^{1/2}] = 2 \tr [\m{A} + \m{B}- \sqrt{\m{Q}_{\varepsilon}}].
\end{align*}
\end{proof}

\end{document}
